\documentclass[UTF8]{ctexart}
\usepackage[a4paper,margin=2.6cm]{geometry}
\usepackage{amsmath, amssymb, amsthm}
\usepackage{graphicx}
\usepackage[colorlinks=true,linkcolor=blue]{hyperref}

\title{复变函数学习笔记：5.1节 Jensen公式}
\author{}
\date{}

\newtheorem{theorem}{定理}[section]
\newtheorem{lemma}[theorem]{引理}
\newtheorem{definition}[theorem]{定义}
\newtheorem{corollary}[theorem]{推论}
\newtheorem{example}{例}[section]
\newtheorem{remark}{注记}[section]
\numberwithin{equation}{section}

\newcommand{\RR}{\mathbb{R}}
\newcommand{\CC}{\mathbb{C}}

\begin{document}

\maketitle

\setcounter{section}{5}
\setcounter{subsection}{0}
\subsection{Jensen公式}

在本节中，我们记 $D_R$ 为以原点为中心、半径为 $R$ 的开圆盘，$C_R$ 为其边界圆周。本节始终排除函数恒为零的平凡情形。Jensen公式建立了全纯函数在圆盘内的零点个数、圆心的模长以及圆周上的对数平均值三者之间的深刻联系，它是研究整函数增长阶与零点分布关系的核心工具。

\begin{theorem}[Jensen公式]
设 $\Omega$ 是包含闭圆盘 $\overline{D_R}$ 的一个开集，$f$ 在 $\Omega$ 内全纯，且 $f(0)\neq0$，$f$ 在圆周 $C_R$ 上无零点。若 $z_1,\dots,z_N$ 是 $f$ 在圆盘 $D_R$ 内的全部零点（按重数计数），则
\begin{equation}\label{eq:jensen}
\log|f(0)| = \sum_{k=1}^{N} \log\left(\frac{|z_k|}{R}\right) + \frac{1}{2\pi} \int_0^{2\pi} \log|f(Re^{i\theta})| \, d\theta .
\end{equation}
\end{theorem}

\begin{proof}[证明思路]
定理的证明分四步进行：
\begin{enumerate}
    \item 验证公式对函数的乘积具有可加性，从而可将一般情形归结为两类特殊函数：无处为零的函数和形如 $z-w$ 的一次因式。
    \item 对于在 $\overline{D_R}$ 内无处为零的函数，可将其写为 $e^{h(z)}$，此时 $\log|f| = \operatorname{Re}(h)$，利用调和函数的均值性质即得结论。
    \item 对于 $f(z)=z-w$（$w\in D_R$），通过计算积分直接验证公式成立。
    \item 将一般函数分解为上述两种函数的乘积，由第1步的可加性完成证明。
\end{enumerate}
\end{proof}

\begin{proof}[详细证明]
\textbf{第1步：乘积的可加性。}
设 $f_1$ 和 $f_2$ 都满足定理条件且公式对它们成立。则对乘积 $f_1f_2$，其零点集是两者零点集的并（按重数相加），而 $\log|f_1f_2| = \log|f_1| + \log|f_2|$。因此公式右边各项分别相加后仍成立。这说明只需对“基本构件”证明公式即可。

\textbf{第2步：无处为零的函数。}
设 $g$ 在 $\overline{D_R}$ 内全纯且无零点。我们断言：
\[
\log|g(0)| = \frac{1}{2\pi} \int_0^{2\pi} \log|g(Re^{i\theta})| \, d\theta .
\]
在比 $D_R$ 稍大的圆盘内，$g$ 仍无零点。由于圆盘是单连通的，根据定理6.2（第3章），存在全纯函数 $h$ 使得 $g(z) = e^{h(z)}$。于是
\[
|g(z)| = |e^{h(z)}| = e^{\operatorname{Re}(h(z))}, \quad \log|g(z)| = \operatorname{Re}(h(z)).
\]
调和函数 $\operatorname{Re}(h)$ 满足均值性质（第3章推论7.3），因此
\[
\operatorname{Re}(h(0)) = \frac{1}{2\pi} \int_0^{2\pi} \operatorname{Re}(h(Re^{i\theta})) \, d\theta ,
\]
这正是所需的等式。

\textbf{第3步：一次因式 $f(z)=z-w$。}
设 $w\in D_R$，需证
\[
\log|w| = \log\left(\frac{|w|}{R}\right) + \frac{1}{2\pi} \int_0^{2\pi} \log|Re^{i\theta} - w| \, d\theta .
\]
由于 $\log(|w|/R) = \log|w| - \log R$ 且
\[
\log|Re^{i\theta} - w| = \log R + \log|e^{i\theta} - w/R|,
\]
代入后两边 $\log R$ 抵消，原等式等价于
\[
\int_0^{2\pi} \log|e^{i\theta} - a| \, d\theta = 0, \quad |a|<1,
\]
其中 $a = w/R$。作变量代换 $\theta\mapsto -\theta$，只需证
\[
\int_0^{2\pi} \log|1 - a e^{i\theta}| \, d\theta = 0, \quad |a|<1.
\]
考虑函数 $F(z) = 1 - az$，它在单位闭圆盘内无零点。如第2步，存在全纯函数 $G$ 使得 $F(z)=e^{G(z)}$，从而 $\log|F(z)| = \operatorname{Re}(G(z))$。由于 $F(0)=1$，有 $\log|F(0)|=0$。对调和函数 $\log|F(z)|$ 应用均值性质，即得所需积分恒等式。

\textbf{第4步：一般情形。}
将 $f(z)$ 在 $\overline{D_R}$ 内的零点 $z_1,\dots,z_N$ 全部提取出来，令
\[
g(z) = \frac{f(z)}{(z-z_1)\cdots(z-z_N)} .
\]
每个 $z_k$ 都是 $g$ 的可去奇点，且补充定义后 $g$ 在 $\overline{D_R}$ 内全纯且无零点。于是 $f(z) = g(z)\prod_{k=1}^N (z-z_k)$。对 $g$ 应用第2步的公式，对每个因式 $(z-z_k)$ 应用第3步的公式，再由第1步的可加性，即得Jensen公式~\eqref{eq:jensen}。
\end{proof}

\subsubsection{Jensen公式的积分形式}

从Jensen公式可以导出一个便于估计的积分恒等式。设 $\mathfrak{n}(r)$ 表示 $f$ 在圆盘 $D_r$ 内的零点个数（按重数计数）。若 $f(0)\neq0$ 且 $f$ 在 $C_R$ 上无零点，则
\begin{equation}\label{eq:jensen-integral}
\int_0^R \mathfrak{n}(r) \frac{dr}{r} = \frac{1}{2\pi} \int_0^{2\pi} \log|f(Re^{i\theta})| \, d\theta - \log|f(0)|.
\end{equation}

\begin{lemma}\label{lem:zero-sum}
若 $z_1,\dots,z_N$ 是 $f$ 在 $D_R$ 内的零点，则
\[
\sum_{k=1}^N \log\left|\frac{R}{z_k}\right| = \int_0^R \mathfrak{n}(r) \frac{dr}{r}.
\]
\end{lemma}
\begin{proof}
首先注意到
\[
\sum_{k=1}^N \log\left|\frac{R}{z_k}\right| = \sum_{k=1}^N \int_{|z_k|}^R \frac{dr}{r}.
\]
定义特征函数 $\eta_k(r) = 1$ 若 $r > |z_k|$，否则为 $0$。则 $\sum_{k=1}^N \eta_k(r) = \mathfrak{n}(r)$。于是
\[
\sum_{k=1}^N \int_{|z_k|}^R \frac{dr}{r} = \sum_{k=1}^N \int_0^R \eta_k(r)\frac{dr}{r} = \int_0^R \mathfrak{n}(r)\frac{dr}{r},
\]
得证。
\end{proof}

将引理代入Jensen公式~\eqref{eq:jensen}，立刻得到积分形式~\eqref{eq:jensen-integral}。这一形式在估计整函数的零点密度时极为方便。

\subsubsection{例子：$\sin \pi z$ 的Jensen公式验证}

\begin{example}
考虑整函数 $f(z) = \sin \pi z$。它在圆盘 $D_R$ 内的零点为整数 $n$ 满足 $|n|<R$。取 $R$ 不是整数，以保证圆周上无零点。由Euler公式，
\[
|\sin \pi z| = \left| \frac{e^{i\pi z} - e^{-i\pi z}}{2i} \right| \le e^{\pi |z|},
\]
故当 $R$ 很大时，$\log|f(Re^{i\theta})| \le \pi R$。另一方面，可以更精确地计算出
\[
\frac{1}{2\pi} \int_0^{2\pi} \log|\sin(\pi R e^{i\theta})| \, d\theta = \log\left(\frac{e^{\pi R}}{2}\right) + o(1) \quad (R\to\infty).
\]
而 $f(0)=0$，Jensen公式需稍作修改（考虑 $f(z)/z$ 等）。但积分形式~\eqref{eq:jensen-integral} 仍能给出零点个数 $\mathfrak{n}(R) \sim 2R$，这与 $\sin \pi z$ 在 $[-R,R]$ 内约有 $2R$ 个零点的事实相符。
\end{example}

\subsubsection{逻辑脉络总结}

\begin{itemize}
    \item \textbf{Jensen公式} 将全纯函数在圆心的模长、圆周上的对数平均值、以及圆内零点的分布三者紧密联系起来。
    \item 其证明通过“分解为基本构件”的策略，充分利用了全纯函数的局部对数存在性和调和函数的均值性质。
    \item \textbf{积分形式} 是研究整函数增长阶与零点分布关系的核心工具，为后续Hadamard因子分解定理奠定了基础。
\end{itemize}

\end{document}